Symbolic objects in Orbiter represent algebraic expressions with coefficients in finite fields.
They consist of the elementary arithmetic operations, 
numbers, variables and exponents.
In Orbiter, an expression can be parsed from a text representation. 
The purpose of parting is to produce an internat representation of the expression as a syntax tree. 
The syntax tree is an efficient way to represent the expression, 
and algebraic operations will be performed by manipulating the syntax tree.
The process of creating the abstract syntax tree 
from the text representation is called parsing. 
Orbiter can deal with vectors and matrices of symbolic expressions. 
These are simply sequences of symbolic objects.


\bigskip


Tables~\ref{tab:symbolic:objects:1}-\ref{tab:symbolic:objects:2} 
show Orbiter commands to create symbolic objects. 
\orbitertableofcommands{symbolic_object_1}
\orbitertableofcommands{symbolic_object_2}

\bigskip

Table~\ref{tab:symbolic:objects:activity} 
show Orbiter activities for symbolic objects. 
\orbitertableofcommands{symbolic_object_activity}

\bigskip


Let us consider some examples. 
The elements of the finite field $\bbF_q$ are coded as integers in the interval from $0$ to $q-1$.  
For more information about finite fields, 
see Chapter~\ref{chapter:finite:fields}.
It is fair to assume that $0$ represents zero and $1$ represents one.
Numbers outside the interval $0$ to $q-1$ are not allowed 
and lead to error messages when encountered.
However, when parsing a algebraic formula from text, 
a negative sign is interpreted as a uniary or binary minus, 
as hence one can say that negative numbers are allowed at the time of input.

\bigskip

The first example creates the polynomial $1-X-X^2$ over $\bbF_7$:

\sourcecode{symbolic_poly}

The output is 
$$
M = 
\begin{array}{*{1}c}
1 + (6 X) + (6 X^{2})\\
\end{array}
$$
The coefficients of $X$ and $X^2$ are $6$, because $6 \equiv -1$ mod $7$ 
and Orbiter chooses field elements of $\bbF_q$ to lie in the interval $[0,q-1].$



\bigskip

The second example creates the polynomial CRC32, 
using the makefile variable

\sourcecodevariable{CRC32_ETHERNET_POLY}

The command 

\sourcecode{symbolic_CRC32}


produces the following output:
$$
M = 
\begin{array}{*{1}c}
1 + X + X^{2} + X^{4} + X^{5} + X^{7} + X^{8} + X^{10} + X^{11} + X^{12} + X^{16} + X^{22} + X^{23} + X^{26} + X^{32}\\
\end{array}
$$



\bigskip


The second example creates the BCH polynomial Alfa, 
defined over the field $\bbF_{256}$, 
using the makefile variable

\sourcecodevariable{BCH_POLYNOMIAL_ALFA_F256}

The command 

\sourcecode{symbolic_CRC_alfa}


produces the following output:
$$
M = 
\begin{array}{*{1}c}
214 + X^{2} + (167 X)\\
\end{array}
$$



\bigskip



The next example creates the CRC polynomial Bravo over the field $\bbF_{256}$.
The makefile variable is:

\sourcecodevariable{BCH_POLYNOMIAL_BRAVO_F256}

The actual command is:

\sourcecode{symbolic_CRC_bravo}

The output is 
$$
M = 
\begin{array}{*{1}c}
1 + X^{4} + (27 X^{2}) + (23 X) + (213 X^{3})\\
\end{array}
$$

\bigskip


%\sourcecode{symbolic_test_4_expand_draw}


Let us consider the formula $(a+b)^3$ and expand:


\sourcecode{symbolic_test_5_expand}


The output is 
$$
M1 = 
\left[
\begin{array}{*{1}c}
a^{3} + b^{3} + (3 a^{2} b) + (3 a b^{2})\\
\end{array}
\right]
$$
It is also possible to inspect the 
representation of the expressions as abstract syntax trees. 
The node labeled R signifies the root node of the tree. 
The tree on the left is the syntax tree before expansion, 
whereas the tree on the right is the syntax tree after expansion:

\orbiteroutput{GRAPHICS/WRAPPERS/symbolic_test_5}

\bigskip


The command 

\sourcecode{test_expand_1}

tests that 
$$
a(x+y) = ax + ay.
$$

\bigskip


The command 

\sourcecode{test_expand_2}

tests that 
$$
a(-(x-y)) = -ax + ay.
$$

\bigskip



The command 

\sourcecode{test_expand_3}

tests that 
$$
a(-(x-y)+(x+y)) = 2ay.
$$
The output is 
$$
fe = 
\begin{array}{*{1}c}
(2 a y)\\
\end{array}
$$


The command shows the abstract syntax trees 
for the intermediate steps in the algebraic simplification.

\bigskip




The command 

\sourcecode{test_expand_4}

tests that 
$$
x*(a+(x*(a+(x*(a+x*(a+x*(a+x*(a+x)))))))) = ax + ax^2+ax^3+ax^4+ax^5+ax^6+x^7.
$$
The output is 
$$
fx = 
\begin{array}{*{1}c}
x^{7} + (a x^{2}) + (a x^{3}) + (a x^{4}) + (a x^{5}) + (a x^{6}) + (a x)\\
\end{array}
$$
The abstract syntax tree before expansion is:
\orbiteroutput{GRAPHICS/WRAPPERS/test_expand_4_before}
The abstract syntax tree after expansion is:
\orbiteroutput{GRAPHICS/WRAPPERS/test_expand_4_after}
The abstract syntax trees for the intermediate steps 
during expansion are avaliable also. 

\bigskip



Objects are stored in the Orbiter symbol table. There is only one table. 
There can be only one object with any given name.
Symbolic objects can recycle objects 
defined earlier using a process called implicit substitution.
In this process, the object being referenced 
is taken from the symbol table and substituted immediately.
For instance, in the next example we create an 
object $a = x + y$ and then another object $b = a+z$ 
which immediately expands into $b = x + y + z.$
This is because the object $a$ already exists. 
Implicit substitution is non-recursive, 
so the substituted syntax tree is not subjected to implicit substitution itself. 
This prevents unlimited recursion which would result in crashes. 


\sourcecode{symbolic_objects_building_on_earlier_objects}

The output of the command is:
$$
b = 
\begin{array}{*{1}c}
x + y + z\\
\end{array}
$$

Self-referential objects are allowed. 
The following command is legal:

\sourcecode{symbolic_objects_self_referential_1}

The output is
$$
x = 
\begin{array}{*{1}c}
x\\
\end{array}
$$

\bigskip

Likewise, the following command is ok:

\sourcecode{symbolic_objects_self_referential_2}

The output is
$$
a = 
\begin{array}{*{1}c}
a + x\\
\end{array}
$$

\bigskip

The implicit substitution process is limited to scalar quantities. 
Matrices and vectors cannot be substituted.




\bigskip



Let us create a Vandermonde matrix. We use the makefile variable 

\sourcecodevariable{VANDERMONDE_4X4_FORMULA}

The following command creates the matrix:

\sourcecode{symbolic_matrix1}

The output is:
$$
\left[
\begin{array}{*{4}c}
1  & x0  & x0^{2}  & x0^{3}\\
1  & x1  & x1^{2}  & x1^{3}\\
1  & x2  & x2^{2}  & x2^{3}\\
1  & x3  & x3^{2}  & x3^{3}\\
\end{array}
\right]
$$

\bigskip

Let us create a fully symbolic $2 \times 2$ matrix:

\sourcecode{symbolic_matrix_2x2}

The output is:
$$
M = 
\left[
\begin{array}{*{2}c}
a  & b\\
c  & d\\
\end{array}
\right]
$$

\bigskip


Let us create the determinant of a symbolic $2 \times 2$ matrix:

\sourcecode{symbolic_determinant_2x2}

The output is
$$
det1 = 
\begin{array}{*{1}c}
(a d) + (12 b c)\\
\end{array}
$$

\bigskip


Let us create the determinant of a symbolic $3 \times 3$ matrix:

\sourcecode{symbolic_determinant_3x3}


The output is
$$
det1 = 
\begin{array}{*{1}c}
(c d h) + (b f g) + (12 b d i) + (12 a f h) + (a e i) + (12 c e g)\\
\end{array}
$$

\bigskip

Let us create the determinant of the Vandermonde matrix from above:


\sourcecode{symbolic_determinant_vandermonde}

The output is lengthy, so it is given in text notation:

{\tt
  0 : \\
(12\ *\ x1\ *\ x2\ \symbol{94}\ 3\ *\ x3\ \symbol{94}\ 2)\ +\ (12\ *\ x1\ \symbol{94}\ 2\ *\ x2\ *\ x3\ \symbol{94}\ 3)\ +\ (x1\ \symbol{94}\ 2\ *\ x2\ \symbol{94}\ 3\\
\ *\ x3)\ +\ (x1\ \symbol{94}\ 3\ *\ x2\ *\ x3\ \symbol{94}\ 2)\ +\ (12\ *\ x1\ \symbol{94}\ 3\ *\ x2\ \symbol{94}\ 2\ *\ x3)\ +\ (12\ *\ x0\ *\ x\\
2\ \symbol{94}\ 2\ *\ x3\ \symbol{94}\ 3)\ +\ (x0\ *\ x2\ \symbol{94}\ 3\ *\ x3\ \symbol{94}\ 2)\ +\ (x0\ *\ x1\ \symbol{94}\ 2\ *\ x3\ \symbol{94}\ 3)\ +\ (12\ *\ x0\\
\ *\ x1\ \symbol{94}\ 2\ *\ x2\ \symbol{94}\ 3)\ +\ (12\ *\ x0\ *\ x1\ \symbol{94}\ 3\ *\ x3\ \symbol{94}\ 2)\ +\ (x0\ *\ x1\ \symbol{94}\ 3\ *\ x2\ \symbol{94}\ 2)\ +\\
\ (x0\ \symbol{94}\ 2\ *\ x2\ *\ x3\ \symbol{94}\ 3)\ +\ (12\ *\ x0\ \symbol{94}\ 2\ *\ x2\ \symbol{94}\ 3\ *\ x3)\ +\ (12\ *\ x0\ \symbol{94}\ 2\ *\ x1\ *\ \\
x3\ \symbol{94}\ 3)\ +\ (x0\ \symbol{94}\ 2\ *\ x1\ *\ x2\ \symbol{94}\ 3)\ +\ (x0\ \symbol{94}\ 2\ *\ x1\ \symbol{94}\ 3\ *\ x3)\ +\ (12\ *\ x0\ \symbol{94}\ 2\ *\ x\\
1\ \symbol{94}\ 3\ *\ x2)\ +\ (12\ *\ x0\ \symbol{94}\ 3\ *\ x2\ *\ x3\ \symbol{94}\ 2)\ +\ (x0\ \symbol{94}\ 3\ *\ x2\ \symbol{94}\ 2\ *\ x3)\ +\ (x0\ \symbol{94}\ 3\\
\ *\ x1\ *\ x3\ \symbol{94}\ 2)\ +\ (12\ *\ x0\ \symbol{94}\ 3\ *\ x1\ *\ x2\ \symbol{94}\ 2)\ +\ (12\ *\ x0\ \symbol{94}\ 3\ *\ x1\ \symbol{94}\ 2\ *\ x3)\ \\
+\ (x0\ \symbol{94}\ 3\ *\ x1\ \symbol{94}\ 2\ *\ x2)\ +\ (x1\ *\ x2\ \symbol{94}\ 2\ *\ x3\ \symbol{94}\ 3)\\
}

\bigskip

The next command computes the characteristic polynomial 
of a symbolic $2 \times 2$ matrix. 
When all is done, the terms of the polynomial 
are collected to form a coefficient vector:


\sourcecode{symbolic_characteristic_polynomial_2x2}

The output is:
$$
coeffs = 
\left[
\begin{array}{*{1}c}
(a d) + (12 b c)\\
(12 d) + (12 a)\\
1\\
\end{array}
\right]
$$


%\sourcecode{symbolic_characteristic_polynomial_2x2_draw}

\bigskip

Let us do the characteristic polynomial of a symbolic $3 \times 3$ matrix:




\sourcecode{symbolic_characteristic_polynomial_3x3}

The output is:
$$
coeffs = 
\left[
\begin{array}{*{1}c}
(c d h) + (b f g) + (12 b d i) + (12 a f h) + (a e i) + (12 c e g)\\
(12 e i) + (c g) + (b d) + (12 a i) + (12 a e) + (f h)\\
a + e + i\\
12\\
\end{array}
\right]
$$


\bigskip

Let us compare Orbiter and Maple with an example. 
We create an object and then substitute values
over the field $\bbF_{13}$. 
Here is the Orbiter code:


\sourcecode{symbolic_object1}

The output is:

$$
M = 
\begin{array}{*{1}c}
a + (10 ((3 d) + (12 a^{2} b^{2} c^{2}) + (2 c^{4})^{3})) + (2 b^{2} (x + (12 y)))\\
\end{array}
$$


\bigskip



Let us substitute some values and evaluate the formula:


\sourcecode{symbolic_object1_evaluate}

The value obtained after substitution is:
$$
eval = 
\left[
\begin{array}{*{1}c}
4\\
\end{array}
\right]
$$


\bigskip


Let us perform the same task in Maple:

\orbiteroutput{GRAPHICS/WRAPPERS/symbolic_maple_1}


\bigskip



%\sourcecode{symbolic_object1_expand}

%\sourcecode{symbolic_object1_draw}

Let us create the general equation 
of the Eckardt surface using a makefile variable:

\sourcecodevariable{ECKARDT_SURFACE_EQN}


The following command parses the equation:


\sourcecode{symbolic_Eckardt_surface}


The output is:
$$
M = 
\begin{array}{*{1}c}
(12 X3 a b^{2} (X0^{2} + X1^{2} + X2^{2})) + (X0 X1 X2 b^{3} (1 + a^{2})) + (X3^{3} a)\\
\end{array}
$$

We specify $a=3$ and $b=1$ by issuing the following command:

\sourcecode{symbolic_Eckardt_surface_q13_a3_b1}

The output is:
$$
M2 = 
\left[
\begin{array}{*{1}c}
(10 X2^{2} X3) + (10 X1^{2} X3) + (10 X0 X1 X2) + (10 X0^{2} X3) + (3 X3^{3})\\
\end{array}
\right]
$$


\bigskip


Let us consider the 4 parameter normal form 
of a cubic surface from~\cite{BettenKaraoglu2022}.
A makefile variable can be used to define the equation in text form:

\sourcecodevariable{F_abcd_eqn}

The following command parses the equation.

\sourcecode{symbolic_Fabcd}

The output in text form is:
{\tt
  0 : \\
(X0\ *\ X1\ *\ X2\ *\ (a\ +\ b\ +\ (12\ *\ c)\ +\ (12\ *\ d))\ *\ ((12\ *\ a\ *\ b\ *\ d)\ +\ (12\ *\ a\ \\
*\ c\ *\ d)\ +\ (b\ *\ c\ *\ d)\ +\ (a\ *\ d)\ +\ (12\ *\ b\ *\ c)\ +\ (a\ *\ b\ *\ c)))\ +\ (X0\ *\ X1\ *\\
\ X3\ *\ (b\ +\ (12\ *\ d))\ *\ ((12\ *\ a\ \symbol{94}\ 2\ *\ d)\ +\ (12\ *\ a\ *\ c\ \symbol{94}\ 2)\ +\ (b\ *\ c\ \symbol{94}\ 2)\ +\ \\
(a\ *\ d)\ +\ (12\ *\ b\ *\ c)\ +\ (a\ \symbol{94}\ 2\ *\ c)))\ +\ (12\ *\ X0\ *\ X2\ \symbol{94}\ 2\ *\ ((12\ *\ a\ *\ b\ *\ \\
d)\ +\ (12\ *\ a\ *\ c\ *\ d)\ +\ (b\ *\ c\ *\ d)\ +\ (a\ *\ d)\ +\ (12\ *\ b\ *\ c)\ +\ (a\ *\ b\ *\ c))\ \\
*\ ((12\ *\ b\ *\ c)\ +\ (a\ *\ d)))\ +\ (12\ *\ X0\ *\ X2\ *\ X3\ *\ (b\ +\ (12\ *\ d))\ *\ ((12\ *\ a\\
\ *\ b\ *\ c\ \symbol{94}\ 2)\ +\ (12\ *\ a\ \symbol{94}\ 2\ *\ d)\ +\ (a\ *\ b\ *\ d)\ +\ (b\ *\ c\ \symbol{94}\ 2)\ +\ (12\ *\ b\ *\ c\ *\\
\ d)\ +\ (a\ \symbol{94}\ 2\ *\ c\ *\ d)))\ +\ (12\ *\ X1\ \symbol{94}\ 2\ *\ X2\ *\ ((12\ *\ a\ *\ b\ *\ d)\ +\ (12\ *\ a\ *\ \\
c\ *\ d)\ +\ (b\ *\ c\ *\ d)\ +\ (a\ *\ d)\ +\ (12\ *\ b\ *\ c)\ +\ (a\ *\ b\ *\ c))\ *\ (a\ +\ (12\ *\ c)\\
))\ +\ (12\ *\ X1\ \symbol{94}\ 2\ *\ X3\ *\ ((12\ *\ a\ *\ b\ *\ d)\ +\ (12\ *\ a\ *\ c\ *\ d)\ +\ (b\ *\ c\ *\ d)\ \\
+\ (a\ *\ d)\ +\ (12\ *\ b\ *\ c)\ +\ (a\ *\ b\ *\ c))\ *\ (a\ +\ (12\ *\ c)))\ +\ (X1\ *\ X2\ \symbol{94}\ 2\ *\ (\\
(12\ *\ a\ *\ b\ *\ d)\ +\ (12\ *\ a\ *\ c\ *\ d)\ +\ (b\ *\ c\ *\ d)\ +\ (a\ *\ d)\ +\ (12\ *\ b\ *\ c)\ +\\
\ (a\ *\ b\ *\ c))\ *\ ((12\ *\ b\ *\ c)\ +\ (a\ *\ d)))\ +\ (X1\ *\ X2\ *\ X3\ *\ ((12\ *\ a\ \symbol{94}\ 2\ *\ b\\
\ *\ d\ \symbol{94}\ 2)\ +\ (11\ *\ a\ \symbol{94}\ 2\ *\ c\ *\ d\ \symbol{94}\ 2)\ +\ (11\ *\ a\ *\ b\ \symbol{94}\ 2\ *\ c\ \symbol{94}\ 2)\ +\ (a\ *\ b\ \symbol{94}\ 2\\
\ *\ c\ *\ d)\ +\ (2\ *\ a\ *\ b\ *\ c\ \symbol{94}\ 2\ *\ d)\ +\ (a\ *\ b\ *\ c\ *\ d\ \symbol{94}\ 2)\ +\ (12\ *\ b\ \symbol{94}\ 2\ *\ c\ \\
\symbol{94}\ 2\ *\ d)\ +\ (12\ *\ a\ \symbol{94}\ 2\ *\ b\ *\ c)\ +\ (a\ \symbol{94}\ 2\ *\ c\ *\ d)\ +\ (a\ \symbol{94}\ 2\ *\ d\ \symbol{94}\ 2)\ +\ (a\ *\ b\\
\ \symbol{94}\ 2\ *\ c)\ +\ (a\ *\ b\ *\ c\ \symbol{94}\ 2)\ +\ (9\ *\ a\ *\ b\ *\ c\ *\ d)\ +\ (12\ *\ a\ *\ c\ \symbol{94}\ 2\ *\ d)\ +\ (\\
a\ *\ c\ *\ d\ \symbol{94}\ 2)\ +\ (b\ \symbol{94}\ 2\ *\ c\ \symbol{94}\ 2)\ +\ (2\ *\ a\ \symbol{94}\ 2\ *\ b\ *\ c\ *\ d)))\ +\ (X1\ *\ X3\ \symbol{94}\ 2\ \\
*\ a\ *\ c\ *\ (b\ +\ c\ +\ (12\ *\ a)\ +\ (12\ *\ b\ *\ c)\ +\ (12\ *\ d)\ +\ (a\ *\ d))\ *\ (b\ +\ (12\ \\
*\ d)))\ +\ (12\ *\ X0\ \symbol{94}\ 2\ *\ X2\ *\ (b\ +\ (12\ *\ d))\ *\ ((12\ *\ a\ *\ b\ *\ d)\ +\ (12\ *\ a\ *\ \\
c\ *\ d)\ +\ (b\ *\ c\ *\ d)\ +\ (a\ *\ d)\ +\ (12\ *\ b\ *\ c)\ +\ (a\ *\ b\ *\ c)))\\
}
\bigskip

The next command expands the equation. 
The domain of coefficients is the finite field $\bbF_{31}$.

\sourcecode{symbolic_Fabcd_expand}

The expanded equation has 127 terms. 


\bigskip

Over the field $\bbF_{31}$,
let us substitute $(a,b,c,d) = (25,5,5,25)$ 
in the general equation.
The following command does the job:

\sourcecode{symbolic_Fabcd_q31_E18}

The output is:
$$
M2 = 
\left[
\begin{array}{*{1}c}
(13 X1 X2 X3) + (22 X1 X2^{2}) + (22 X1^{2} X3) + (22 X1^{2} X2) + (9 X0 X2^{2}) + (9 X0^{2} X2) + (22 X1 X3^{2})\\
\end{array}
\right]
$$


\bigskip

It is also possible to use implicit 
substitution to achieve the same effect, 
as shown in the following command:

\sourcecode{symbolic_Fabcd_q31_E18_with_implicit_substitution}


The output is again 
$$
M2 = 
\begin{array}{*{1}c}
(13 X1 X2 X3) + (22 X1 X2^{2}) + (22 X1^{2} X3) + (22 X1^{2} X2) + (9 X0 X2^{2}) + (9 X0^{2} X2) + (22 X1 X3^{2})\\
\end{array}
$$

We can use the option \verb'-do_not_simplify' 
to see the syntax tree 
exactly after the moment of substitution, 
before the nodes are evaluated:
So, the command 


\sourcecode{symbolic_Fabcd_q31_E18_with_implicit_substitution_do_not_simplify}

produces the output

{\tt
  0 : \\
((30\ *\ ((25\ *\ 5\ *\ 5)\ +\ (30\ *\ (25\ *\ 5\ *\ 25))\ +\ (30\ *\ (25\ *\ 5\ *\ 25))\ +\ (5\ *\ 5\ \\
*\ 25)\ +\ (25\ *\ 25)\ +\ (30\ *\ (5\ *\ 5))))\ *\ (5\ +\ (30\ *\ 25))\ *\ X0\ \symbol{94}\ 2\ *\ X2)\ +\ (((2\\
5\ *\ 5\ *\ 5)\ +\ (30\ *\ (25\ *\ 5\ *\ 25))\ +\ (30\ *\ (25\ *\ 5\ *\ 25))\ +\ (5\ *\ 5\ *\ 25)\ +\ (2\\
5\ *\ 25)\ +\ (30\ *\ (5\ *\ 5)))\ *\ (25\ +\ 5\ +\ (30\ *\ 5)\ +\ (30\ *\ 25))\ *\ X0\ *\ X1\ *\ X2)\ \\
+\ (((25\ \symbol{94}\ 2\ *\ 5)\ +\ (30\ *\ (25\ \symbol{94}\ 2\ *\ 25))\ +\ (30\ *\ (25\ *\ 5\ \symbol{94}\ 2))\ +\ (5\ *\ 5\ \symbol{94}\ 2)\ \\
+\ (25\ *\ 25)\ +\ (30\ *\ (5\ *\ 5)))\ *\ (5\ +\ (30\ *\ 25))\ *\ X0\ *\ X1\ *\ X3)\ +\ (30\ *\ (((2\\
5\ *\ 25)\ +\ (30\ *\ (5\ *\ 5)))\ *\ ((25\ *\ 5\ *\ 5)\ +\ (30\ *\ (25\ *\ 5\ *\ 25))\ +\ (30\ *\ (25\\
\ *\ 5\ *\ 25))\ +\ (5\ *\ 5\ *\ 25)\ +\ (25\ *\ 25)\ +\ (30\ *\ (5\ *\ 5)))\ *\ X0\ *\ X2\ \symbol{94}\ 2))\ +\ (\\
30\ *\ (((25\ \symbol{94}\ 2\ *\ 5\ *\ 25)\ +\ (30\ *\ (25\ *\ 5\ *\ 5\ \symbol{94}\ 2))\ +\ (30\ *\ (25\ \symbol{94}\ 2\ *\ 25))\ +\ \\
(25\ *\ 5\ *\ 25)\ +\ (5\ *\ 5\ \symbol{94}\ 2)\ +\ (30\ *\ (5\ *\ 5\ *\ 25)))\ *\ (5\ +\ (30\ *\ 25))\ *\ X0\ *\ \\
X2\ *\ X3))\ +\ (30\ *\ ((25\ +\ (30\ *\ 5))\ *\ ((25\ *\ 5\ *\ 5)\ +\ (30\ *\ (25\ *\ 5\ *\ 25))\ +\ \\
(30\ *\ (25\ *\ 5\ *\ 25))\ +\ (5\ *\ 5\ *\ 25)\ +\ (25\ *\ 25)\ +\ (30\ *\ (5\ *\ 5)))\ *\ X1\ \symbol{94}\ 2\ *\\
\ X2))\ +\ (30\ *\ ((25\ +\ (30\ *\ 5))\ *\ ((25\ *\ 5\ *\ 5)\ +\ (30\ *\ (25\ *\ 5\ *\ 25))\ +\ (30\ \\
*\ (25\ *\ 5\ *\ 25))\ +\ (5\ *\ 5\ *\ 25)\ +\ (25\ *\ 25)\ +\ (30\ *\ (5\ *\ 5)))\ *\ X1\ \symbol{94}\ 2\ *\ X3)\\
)\ +\ (((25\ *\ 25)\ +\ (30\ *\ (5\ *\ 5)))\ *\ ((25\ *\ 5\ *\ 5)\ +\ (30\ *\ (25\ *\ 5\ *\ 25))\ +\ (\\
30\ *\ (25\ *\ 5\ *\ 25))\ +\ (5\ *\ 5\ *\ 25)\ +\ (25\ *\ 25)\ +\ (30\ *\ (5\ *\ 5)))\ *\ X1\ *\ X2\ \symbol{94}\\
\ 2)\ +\ ((((1\ +\ 1)\ *\ 25\ \symbol{94}\ 2\ *\ 5\ *\ 5\ *\ 25)\ +\ (30\ *\ (25\ \symbol{94}\ 2\ *\ 5\ *\ 25\ \symbol{94}\ 2))\ +\ (30\\
\ *\ ((1\ +\ 1)\ *\ 25\ \symbol{94}\ 2\ *\ 5\ *\ 25\ \symbol{94}\ 2))\ +\ (30\ *\ ((1\ +\ 1)\ *\ 25\ *\ 5\ \symbol{94}\ 2\ *\ 5\ \symbol{94}\ 2))\ \\
+\ (25\ *\ 5\ \symbol{94}\ 2\ *\ 5\ *\ 25)\ +\ ((1\ +\ 1)\ *\ 25\ *\ 5\ *\ 5\ \symbol{94}\ 2\ *\ 25)\ +\ (25\ *\ 5\ *\ 5\ *\ 25\\
\ \symbol{94}\ 2)\ +\ (30\ *\ (5\ \symbol{94}\ 2\ *\ 5\ \symbol{94}\ 2\ *\ 25))\ +\ (30\ *\ (25\ \symbol{94}\ 2\ *\ 5\ *\ 5))\ +\ (25\ \symbol{94}\ 2\ *\ 5\ \\
*\ 25)\ +\ (25\ \symbol{94}\ 2\ *\ 25\ \symbol{94}\ 2)\ +\ (25\ *\ 5\ \symbol{94}\ 2\ *\ 5)\ +\ (25\ *\ 5\ *\ 5\ \symbol{94}\ 2)\ +\ (30\ *\ ((1\ \\
+\ 1\ +\ 1\ +\ 1)\ *\ 25\ *\ 5\ *\ 5\ *\ 25))\ +\ (30\ *\ (25\ *\ 5\ \symbol{94}\ 2\ *\ 25))\ +\ (25\ *\ 5\ *\ 25\ \symbol{94}\\
\ 2)\ +\ (5\ \symbol{94}\ 2\ *\ 5\ \symbol{94}\ 2))\ *\ X1\ *\ X2\ *\ X3)\ +\ (5\ *\ 25\ *\ ((25\ *\ 25)\ +\ (30\ *\ (5\ *\ 5\\
))\ +\ (30\ *\ 25)\ +\ 5\ +\ 5\ +\ (30\ *\ 25))\ *\ (5\ +\ (30\ *\ 25))\ *\ X1\ *\ X3\ \symbol{94}\ 2)\\
}



%\sourcecode{symbolic_matrix_minor}

%\sourcecode{symbolic_matrix_nullspace_1}

%\sourcecode{symbolic_matrix_nullspace_2}


%\sourcecode{symbolic_Clebsch}

%\sourcecode{Clebsch_map_up_Fabcd}





\bigskip

